% LaTeX settings shared by the notes and the sheets (PDF output only).
% amsmath, amssymb, graphicx, xcolor and hyperref are already loaded by Quarto.
% HTML equivalents of the colours and boxes are in style.css; keep them in sync.

% Module colours (the HTML ones are in theme-light.scss / theme-dark.scss)
\definecolor{accent}{HTML}{440379}   % Warwick aubergine
\definecolor{alert}{HTML}{B00020}

\usepackage{tcolorbox}
\tcbuselibrary{most}

% ::: {.infobox}
\newtcolorbox{infobox}{enhanced, colback=accent!4, colframe=accent!60,
  boxrule=0.5pt, arc=2pt, left=8pt, right=8pt, top=6pt, bottom=6pt,
  before skip=6pt, after skip=14pt, before upper={\setlength{\parskip}{6pt}}}

% [2]{.marks}: marks pushed to the right margin (moves to the next line if needed)
\newcommand{\marksbox}[1]{\unskip\nobreak\hfil\penalty50\hskip1em\hbox{}\nobreak
  \hfill\mbox{\bfseries\color{accent}#1}}

% [text]{.marking} and ::: {.marking}: notes for markers (staff builds only)
\newcommand{\markingfont}{\small\itshape\color{accent!85!black}}
\newtcolorbox{markingbox}{enhanced, frame hidden, sharp corners,
  colback=accent!4, borderline west={2pt}{0pt}{accent!60},
  left=6pt, right=6pt, top=3pt, bottom=3pt, before skip=4pt, after skip=6pt,
  fontupper=\markingfont, before upper={\setlength{\parskip}{4pt}}}

% [2]{.difficulty}: icons (style set by difficulty-icon in _module.yml)
\newcommand{\difficultystar}{\textcolor{accent}{$\bigstar$}}
\newcommand{\difficultynumber}[1]{\textcolor{accent}{(#1)}}
% Spring icon; the coordinates are generated by module.lua
% (24 units tall = 1em, sitting slightly below the baseline like an emoji)
\usepackage{tikz}
\DeclareRobustCommand{\springicon}[1]{\begin{tikzpicture}[x=0.0417em, y=0.0417em,
  baseline=0.125em]\draw[accent, line width=0.5pt, line join=round,
  line cap=round] plot coordinates {#1};\end{tikzpicture}}
